% !TeX root = ../drafts/e-recursion.tex
\section{Recursion}\label{annex:rec}

A recursion proof shows that leanVM proofs verify. It is not a run of the RISC-V machine: a second machine, the \emph{recursion machine}, proves that a fixed list of rows is satisfied, with the same bus (\S\ref{sec:gp}, \S\ref{sec:gkr}), table sumcheck (\S\ref{sec:air}), Flock (Annex~\ref{annex:flock}) and stacked opening (\S\ref{sec:stacking}, Annex~\ref{annex:pcs}) as a leanVM proof. This annex describes that machine: its wires (\S\ref{rec:wires}), its tables (\S\ref{rec:tables}), its statement (\S\ref{rec:statement}) and its proof (\S\ref{rec:proof}), then the rows that verify a leanVM proof (\S\ref{rec:verifier}), and the trees that aggregate many proofs into one (\S\ref{rec:tree} to \S\ref{rec:root}). Writing a verifier as its rows is the recursive composition of~\cite{BCCT13}.

\subsection{Circuits and wires}\label{rec:wires}

A \emph{circuit} $\rcirc$ is a list of rows, each in one of six tables, and a partition of the rows' slots into \emph{wires}. A wire holds one of three kinds of value: a word of $\K$, an element of $\E$ as its three limbs, or a digest as four words. Every row names one wire per slot. The circuit is public: the verifier builds it from what it verifies, never from a proof, by the same code the prover runs with every value read from a proof replaced by zero. So no row exists or is skipped according to a value; only the prover's assignment of values to wires depends on a proof.

The slots naming one wire are the equalities the circuit asserts, and they are checked on the bus by copy cycles, the permutation argument of PLONK~\cite{GWC19} in the multiset form of \S\ref{sec:gp}. Table $j$ has $2^{\tau_j}$ rows, with $\tau_j\le32$, and its slots have indices $\rgid$ among all tables' slots, in table order. Slot $\rgid$ of row $z$ has the key
\[
  \slotkey(\rgid,z)\;=\;\intk{2^{32}\rgid+z},
\]
an index column (\S\ref{sec:idxcol}), and since $z<2^{32}$ no two slots of the circuit share a key. Order the slots of each wire by table, row and slot. The slot carrying value $v$ \pull{}s $\tup{\slotkey,v_0,v_1,v_2,v_3}$ at its own key and \push{}es $\tup{\slotnext,v_0,v_1,v_2,v_3}$, unused limbs zero, where $\slotnext$ is the key of the next slot of its wire, the last slot's next being the first. On a padding row every slot has $\slotnext=\slotkey$. The column $\slotnext$ is public, a function of the circuit.

\begin{lemma}[Copy cycles]\label{lem:rec-cycles}
If the pushed and pulled multisets are equal, every slot of a wire carries the same value.
\end{lemma}
\begin{proof}
The keys are distinct, so each key is the first coordinate of exactly one pulled tuple. The map $\slotnext$ is a permutation of the keys: each wire's slots form one cycle under it, and a padding row's slot is a cycle of one. So each key is also the first coordinate of exactly one pushed tuple, that of the slot whose next it is. Equal multisets therefore match, at every key $k$, the tuple slot $k$ pulls with the tuple its predecessor pushes, and the two carry the same value. Around a cycle all values are one.
\end{proof}

By Theorem~\ref{th:product_check} the bus's product check passes on unequal multisets with probability at most $4\cdot2^{\mu}/|\E|$, $2^{\mu}$ the larger side's number of leaves. Lemma~\ref{lem:gp} needs no injectivity of the fingerprint beyond what it proves: distinct tuples have distinct fingerprint polynomials, and the product check's error is the only probability in the argument.

A wire in one slot is a cycle of one: the slot pulls and pushes the same tuple, which constrains nothing. That is how a value the prover chooses enters: a scalar of the proof, a hint, a bit of a word. A padding row's slots are such cycles too, and no real slot's $\slotnext$ names a padding key, so whatever a padding row carries cancels and touches no wire. An equality the circuit asserts merges two wires into one. A constant and a word of the statement are rows of the \textsf{PUB} table, a framework block (\S\ref{sec:leafstack}) whose every coordinate is public, so a wire with a \textsf{PUB} slot holds the value the verifier gives it. A product by the constant $1$ or $0$, or a sum with the constant $0$, emits no row, and neither does an \textsf{EMUL} or \textsf{EXK} row whose input wires are an earlier row's, the factors of $a\cdot b+d$ in either order and the terms of a sum $1\cdot a+d$ too: its output is the earlier row's wire. The constants are wires of the circuit itself, so this depends on its structure and never on a value.

\subsection{Tables}\label{rec:tables}

A table's slots are degree-two forms over its columns (\S\ref{sec:m3}), so an output need not be committed: it rides the bus as the form that computes it.
\begin{itemize}[itemsep=1pt,topsep=3pt]
\item \textsf{EMUL}: $c=a\cdot b+d$ over $\E$, with the nine columns $a,b,d$. The slot of $c$ carries the product modulo $y^3+y+1$: $c_0=a_0b_0+a_1b_2+a_2b_1+d_0$, $c_1=a_0b_1+a_1b_0+a_1b_2+a_2b_1+a_2b_2+d_1$ and $c_2=a_0b_2+a_1b_1+a_2b_0+a_2b_2+d_2$.
\item \textsf{EXK}: $c=a\cdot k+d$ with $k\in\K$, seven columns, $c_i=a_ik+d_i$.
\item \textsf{HASH}: one BLAKE2s compression. Its words are ports of a packed witness of the \textsf{HASH} class circuit (\S\ref{flock:classes}), which Flock proves, and its one committed column is a bit $b$ with the identity $b^2=b$.
\item \textsf{SPLIT}: a word and its 64 bits, with $w=\sum_i\intk{2^i}\,b_i$ and $b_i^2=b_i$. Read one way it splits a word, read the other it packs bits into one.
\item \textsf{CAST}: four words seen as a digest, as an element of $\E$, as two elements of $\E$ with a zero top limb (a digest's two halves on the stream), and as four words.
\end{itemize}
A hash row's slots are the chaining value $h$, the counter and the finalization word, the message's first half through the mux, $b$, the message words $m_4,m_5,m_6$ as an element of $\E$, $m_7$, the output, the output's first three words as an element of $\E$, and the eight message words alone. The mux carries $m_i+b\,(m_i+m_{4+i})$, the left half when $b=0$ and the right one when $b=1$. A duplex compression (\S\ref{sec:e2e-fiat-shamir}) wires the chaining value, the framed counter and final flag, and every message word; a seed fixes the chaining input to the parameter IV. Absorbed scalars are cast to words and buffered across 64-byte blocks. Output-node digests are cast to four words, retained across samples, and the next three words are cast back to $\E$, even when they span two output blocks. Role, final-length and consumed-output counters are circuit constants determined by the verifier schedule, never unconstrained witness choices. A Merkle node on a query's path wires its child to the mux, its direction to $b$, and leaves its sibling free; a node of a query batch's shared subtree (\S\ref{rec:verifier}) wires its left child to the mux, $b$ to $0$, and its right child to $m_4,\dots,m_7$ through a \textsf{CAST} row. A block of an ordinary long hash wires $h$ to the previous block's output and the counter to the block's byte count. Every hash row is a row of this one table, so one Flock reduction proves them all.

Each table's summand in the table sumcheck is its push form, its pull form weighted $\xi$, and its identities, the $i$-th of table $j$ weighted $\xi^{2+o_j+i}$ with $o_j$ the identities of the tables before it. The powers are distinct, so the target pins each identity's sum to zero.

\subsection{The statement}\label{rec:statement}

The statement is the list of the \textsf{PUB} rows' words that the circuit exposes. The transcript starts from a digest naming the circuit and absorbs, before anything is sent, the BLAKE2s hash of the statement's limbs as little-endian bytes, eight a block from the parameter IV, the last block's counter their length in bytes. Every challenge therefore depends on the statement.

The bus alone does not bind it. The verifier evaluates the \textsf{PUB} block's public coordinates at the bus's point, after $\alpha$, $\beta$ and $\zeta$ are drawn, and a change $\delta$ to a word's limbs moves both sides' totals by the same $\sum_l w_l\delta_l$, with weights $w_l\in\E$ and $\delta_l\in\K$. That sum vanishes on a space of $\K$-linear dimension at least one in the four limbs of a word. A statement moved inside it leaves every check of the bus and the table sumcheck as it was, so it is the seed that refuses it: a different statement is a different transcript.

\subsection{The proof}\label{rec:proof}

A recursion proof is a leanVM proof of the recursion machine's run. The stack (\S\ref{sec:stacking}) holds the \textsf{HASH} table's packed witness and then every committed column, \textsf{PUB} committing none. The bus balances the copy cycles by the GKR of \S\ref{sec:gkr}; the verifier then draws $\xi$, and one table sumcheck proves every table's summand at the bus's point. Flock reduces the hash rows to one ring-switched claim on the packed witness, its matrices settled against the compression circuit, and one opening discharges every column claim and that claim (Annex~\ref{annex:pcs}). A circuit whose witness exceeds one commitment has no proof: the prover refuses it, and so does the verifier.

\subsection{The verifier in rows}\label{rec:verifier}

The rows that verify a leanVM proof are the verifier's core (\S\ref{sec:e2e}): every check that depends on the proof, which leaves claims on polynomials only the program and the VM's circuits fix. Those claims are the rows' output, as wires, and whoever holds the program and the circuits settles them later.

\paragraph{One verifier, two arithmetics.} The core is written once over an arithmetic of $\E$ and a transcript. Natively an element is a value, a read takes the next scalar of the proof, and a failed equality refuses the proof. In rows an element is a wire, a read is a free wire that a hash row absorbs into the transcript, a product $a\cdot b+d$ is an \textsf{EMUL} row and $a\cdot k+d$ with a constant $k\in\K$ an \textsf{EXK} row, and an equality merges two wires, which the copy cycles of Lemma~\ref{lem:rec-cycles} enforce. The bus's GKR (\S\ref{sec:gkr}), its decomposition into column claims (\S\ref{sec:leafstack}), the table sumcheck (\S\ref{sec:air}), whose final value is rebuilt from the bus forms and \texttt{EXT}'s three identities (\S\ref{sec:tab-classes}), and the program's claim are this shared code, so the rows check what the native verifier checks by construction. The Flock reductions and the opening are written for the rows alone, and tests hold them to the native verifier claim for claim.

\paragraph{Why the rows never depend on values.} The arithmetic exposes no value to the code above it, so no branch can follow one: a check is an equality, never a test. Where the rows need a value of their own (a hint, the bits of a word), the prover computes it and the rows constrain it. The verifier builds the same rows from the shape of what is verified (the program, each table's height and the rate) with zeros in place of the proof, so the circuit a prover proves is the one a verifier rebuilds.

\paragraph{The announcement.} Each table's height and the rate, read off the stream, are held to the shape's constants. The final clock $\tsfinal$ is cast to its three limbs, whose two upper ones are held to zero, and the low one is split into bits: bit $40$ is held to one, every higher bit and the five low ones to zero (\S\ref{sec:state}). The final state tuple carries $\tsfinal$ twice; the rows build the bus's layout with the clock zero and add its share $\mathrm{eq}_{\mathrm{ST}}\,(w_2+w_3)\,\tsfinal$ to the pull side's total afterwards, $\mathrm{eq}_{\mathrm{ST}}$ the final state block's selector and $w_2,w_3$ the fingerprint weights of its clock coordinates, since a leaf is affine in each coordinate. The native verifier does the same.

\paragraph{Flock.} The batched zerocheck and lincheck over every packed witness (\S\ref{flock:batch}) are replayed as their transcripts and the arithmetic that derives their claims. Each stage's terminal identity, \eqref{flock:eq:batch-zc-terminal} and \eqref{flock:eq:batch-lc-terminal}, is one equality of wires; the latter holds against the form values $\lcmat{f}$ the prover sends, each of which is circuit $f$'s deferred matrix claim. The univariate skip (\S\ref{flock:zerocheck}) interpolates over $\skipdom\cup\skipcoset$. Its vanishing polynomial $\vansub_{\skipdom}(Y)=\prod_{h\in\skipdom}(Y+h)=\sum_{j\le\kskip}c_jY^{2^j}$ is linearized with coefficients in $\K$, since $\skipdom$ is an $\Ftwo$-subspace of $\K$, so its value at $\zskip$ costs $\kskip$ squarings and $\kskip$ products by constants, and $\vansub_{\skipcoset}(\zskip)=\vansub_{\skipdom}(\zskip)+\vansub_{\skipdom}(\fembed(64))$. With $D_n=\bigl(\prod_{0<i<n}\fembed(i)\bigr)^{-1}$, the first round's message $P|_{\skipcoset}$ at $\zskip$ and the lincheck's slices $s$ against the Lagrange basis of $\skipdom$ are
\[
  P(\zskip)=D_{128}\,\vansub_{\skipdom}(\zskip)\,\vansub_{\skipcoset}(\zskip)\sum_{h\in\skipcoset}\frac{P(h)}{\zskip+h},\qquad
  \sum_i\lag_i(\zskip)\,s_i=D_{64}\,\vansub_{\skipdom}(\zskip)\sum_{i}\frac{s_i}{\zskip+\fembed(i)}.
\]
Each inverse $u=1/(\zskip+h)$ is a hint held to $\zskip\,u+h\,u=1$, three rows a node with the sum. An honest prover fails only when $\zskip\in\skipdom\cup\skipcoset$, with probability $128/|\E|$. One $\zskip$ serves the whole batch, so the inverses at $\skipdom$ are taken once and each circuit's slices cost one product a node; the circuits of one block size share $\eq(\fc_{\mathrm{in},f},\fc'_{\mathrm{in},f})$, and the challenges a circuit sits out are one product of suffixes for all. A recursion proof's own \textsf{HASH} rows are a batch of one circuit.

\paragraph{The opening.} The ring-switched family's target and its weight at the terminal point (\S\ref{rs:family}, \S\ref{rs:weight}) are computed in rows, once per verified proof: the map's Frobenius coefficients from its six challenges, the family's slices $\sum_j\rsgamma^j s_{j,i}$, the target by the linearized Horner rule, and each claim's terms at the point's inverse Frobenius ladders, the claims whose points are prefixes of one another, the same wires, sharing one pass of the products. The WHIR verifier (Annex~\ref{annex:pcs}) is replayed level by level. A query's index is cut from a challenge's limbs, split into bits, and each bit is a Merkle direction: the bit is the \textsf{HASH} row's selector $b$, Boolean by its \textsf{SPLIT} row and by the table's identity $b^2=b$, and the row hashes the child on the left or on the right by the mux, so the rows are the same for every index. The queries are stratified (Definition~\ref{def:strata}): a query's fixed top bits are constants of the shape, and its leaf is hashed up its own bits only, to the node of the tree's level $\strs$ that its fixed bits name, from the low siblings of its path as the raw proof writes them out. The batch's largest group reaches every node of the top subtree, whose internal nodes the rows hash once, and every other query's node is merged with the subtree's, so every query's path is complete and the rows are the same for every index. A level-$0$ leaf image starts with the zeros of the absent lanes; its whole zero blocks are one public chaining value, a constant, and the rows hash only the image's tail. A one-block leaf and a node are the same compression, from the parameter IV at counter $64$; what keeps a node from being opened as a leaf is the path's length, the tree's height, which the shape fixes and never the proof.

\paragraph{What the rows leave.} The core's output is the deferred claims (the program's point and value, and each of the $21$ Flock circuits' matrix form and value, \texttt{EXT} having a clock circuit and no class circuit) and a terminal commitment to the transcript, including its consumed-output cursor, not its raw chaining value. If the rows are satisfied, the free wires that the transcript read are a proof that the native core accepts with exactly these claims, by the equalities above and the hash rows' replay of the Fiat-Shamir transcript; settling the claims against the program and the circuits completes the verifier.

\subsection{Aggregation trees}\label{rec:tree}

An aggregation tree proves that many proofs verify, with one proof a native verifier checks once. Its leaves come in $\nfam\ge1$ \emph{families}: leanVM proofs of one program $P$ of one shape (each table's height and the rate), at most one such family, or recursion proofs of one \emph{leaf circuit} $\rleaf$ at its own heights, rate and transcript seed, such as a circuit verifying a batch of signatures. Its other proofs are recursion proofs of $\nfam+1$ circuits, at a rate of their own. A first-level node's rows are mostly its leaves' Merkle paths, one per query, and a lower rate needs fewer queries for the same soundness, so leaves at a lower rate than the tree's proofs move work from the aggregator, who proves every tree proof, to the leaves' provers, who run in parallel.
\begin{itemize}[itemsep=1pt,topsep=3pt]
\item A \emph{first-level node} of family $i$ is a proof of $\rfirst^{(i)}$, which verifies $\arityz$ leaves of that family in rows (\S\ref{rec:verifier} for leanVM proofs, \S\ref{rec:child} for a leaf circuit's) and reduces the claims they leave.
\item A \emph{node} is a proof of $\rnode$, which verifies $\arity\ge2$ tree proofs of any kind in rows (\S\ref{rec:child}) and reduces the claims they leave and carry.
\end{itemize}
The root is a node, or a first-level node. A tree of one first-level node over one leaf, $\arityz=1$, is the recursion of a single proof. A balanced tree has $\arityz\,\arity^d$ leaves, $d$ the number of levels of nodes; since a node verifies any kind, one tree mixes families and depths.

Every tree proof states the same words, whatever the tree's size:
\begin{itemize}[itemsep=1pt,topsep=3pt]
\item its kind's code $\kbit$: $0$ for a first-level node of family $0$, $1$ for a node, $i+1$ for a first-level node of family $i>0$;
\item its digest $\tdig$ of what the leaves under it state, as two elements of $\E$ with a zero top limb;
\item one claim on each \emph{dense} polynomial the tree has, at prefixes of one point $\fdense$: the program's stacked bytecode table $\bctab$ (\S\ref{sec:e2e-bc}) and RAM's image $\imgpoly$ zero padded to a power of two, if leanVM proofs are leaves; the leaf circuits' fixed polynomial $\fixleaf$, if leaf circuits' proofs are; and the nodes' fixed polynomial $\fixnode$ (\S\ref{rec:identity});
\item for each of the $21$ Flock circuits $f$, the values of its matrices $A_0$ and $B_0$ (\S\ref{flock:native}) at the prefixes of length $k_f$ of one row point $\frow$ and one column point $\fcol$, $2^{k_f}$ the circuit's instance width.
\end{itemize}
The kind is a constant of each circuit, the digest is computed in rows, and the claims are the outputs of the node's reduction (\S\ref{rec:reduce}). A first-level node reduces claims only on the polynomials its family fixes ($\bctab$ and $\imgpoly$, or $\fixleaf$): its other values are zero, and a parent weighs them by zero (\S\ref{rec:reduce}).

\paragraph{Digests.} A first-level node's digest is the BLAKE2s hash, from the parameter IV, of a block holding a tag and $\arityz$ and then what its leaves state in order: a leanVM proof's output, or the hash of a leaf circuit's statement, the hash its transcript is seeded by, the leaf circuit's seed then following the count in the block. A node's is that of another tag, $\arity$ and its children's digests in order. The tags and the seeds separate the levels and the families, and the counts fix the items, so a digest names its whole subtree. The root's verifier is given the expected tree, its shape and each leaf's statement, and recomputes the digest from it.

\subsection{A recursion proof in rows}\label{rec:child}

A node verifies each child with the recursion machine's verifier as rows. The bus's GKR, its decomposition and the table sumcheck are the native verifier's code over rows, as for a leanVM proof (\S\ref{rec:verifier}): each table's summand is its two bus forms and its identities (\S\ref{rec:tables}), evaluated at the claimed columns. The \textsf{HASH} table's Flock reduction and the opening are the replays a leanVM proof's verifier uses. The \textsf{HASH} reduction leaves a claim on the matrices of the \textsf{HASH} class circuit, one of the $21$, which joins its reduction.

What the native verifier reads off the child's circuit are its public bus columns: each slot's $\slotnext$ and the \textsf{PUB} rows' four value columns. The child's statement fills the first \textsf{PUB} rows, so the rows evaluate the statement's share of each value column from the statement's words, $\sum_{z<S}\eq(z,\zeta)\,v_j(z)$, in rows linear in the statement. The rest is fixed by the child's circuit: the $\slotnext$ columns and the \textsf{PUB} constants, its \emph{fixed columns}. Each evaluation the bus needs of a fixed column, at the bus point's prefix of the column's height, is a hint the prover supplies and the rows use, and it becomes a claim on $\fixnode$.

A first-level node over a leaf circuit's proofs verifies each the same way, with the leaf circuit's heights, rate and seed, which are constants of $\rfirst^{(i)}$: the leaf's statement is free words the rows hash into its transcript's seed, that hash is the digest's item, and its hints become claims on $\fixleaf$. Such a leaf carries no claim of its own, so the node reduces its fresh claims only.

\subsection{Circuit identity}\label{rec:identity}

Every kind's circuit has the same tables' heights $\taufix$, each table padded to the largest, so one node circuit verifies a child of any kind with the same rows. Each circuit's fixed columns are stacked, largest first at aligned offsets, into a table of $2^{\fixdim}$ words; the stacks share their layout, and with $b=\lceil\log_2(\nfam+1)\rceil$ the fixed polynomial
\[
  \fixnode(x,\kbit_0,\dots,\kbit_{b-1})\;=\;\text{the stack at $x$ of the circuit of code }\textstyle\sum_j\kbit_j2^j,\ \text{zero past the codes}
\]
has $\fixdim+b$ variables, the last ones the kind's bits. Column $c$ of the child's circuit at the bus point's prefix $\zeta_{<\tau_c}$ is $\fixnode(\zeta_{<\tau_c},\mathsf{blk}_c,\kbit_0,\dots,\kbit_{b-1})$, $\mathsf{blk}_c$ the bits of the column's offset over its height. The child's kind $\kbit$ is a word of its statement; the node's rows hold each bit Boolean by $\kbit_j^2=\kbit_j$ in $\E$ and $\kbit=\sum_j\kbit_jX^j$ (for $b=1$ the bit is $\kbit$), and, when $\nfam+1<2^b$, the sum of the codes' indicators $\prod_j(\kbit_j\text{ or }1+\kbit_j)$ to one, so the code is a kind of the tree. Each circuit holds its own kind word to its constant through the bus, so a child is checked against the fixed columns of the circuit it claims to be, and the claims are settled against the circuits the root's verifier builds from the families and the shape. No commitment to a circuit is needed, and $\rnode$ refers to its own fixed columns only through such claims, so it is well defined. The leaf circuits' fixed polynomial $\fixleaf$ is their stacks the same way, each at its own layout zero padded to the largest, selected by public bits naming the circuit's slot: a first-level node's family is a constant of its circuit.

\paragraph{The shared heights.} A first-level node's circuit depends on $\taufix$ only through its statement's length; $\rnode$ verifies children of heights $\taufix$. The heights are the least fixed point of $\tau\mapsto\max(\tau,\max_i\tau(\rfirst^{(i)}(\tau)),\tau(\rnode(\tau)))$, table by table, from the first-level circuits' own heights. The map is monotone and bounded by the heights one commitment holds, so Kleene iteration reaches the fixed point; a node's rows grow with the logarithm of its children's heights, so it takes few rounds, and the builder refuses circuits that have not converged after eight. A leaf circuit's proofs keep their own heights. Every tree proof's transcript starts from a digest of the families (the program and the leaves' shape, or a leaf circuit's seed, statement length, heights and rate), $\arityz$, $\arity$, $\taufix$, the rate and the statement's length. The digest is the same for every kind: the kind is the statement's first word, so the hash of the statement the transcript absorbs next (\S\ref{rec:statement}) separates the kinds' transcripts.

\subsection{The per-node reduction}\label{rec:reduce}

A first-level node over leanVM proofs leaves, for each leaf, the program claim and each Flock circuit's matrix claim (\S\ref{rec:verifier}). The program claim is $\sum_i\mu_i\,\varphi^i(\bctab(\varphi^{-i}(\chi),\vec\alpha))+w\,\imgpoly(p)=v$, $\varphi$ the Frobenius $x\mapsto x^2$ and $\mu_i$ the weight of multiplicity bit $i$. It becomes point claims: the prover supplies $D_i=\bctab(\varphi^{-i}(\chi),\vec\alpha)$ for each bit and $I=\imgpoly(p_{<m})$, and the rows check $\sum_i\mu_iD_i^{2^i}+w\prod_{j\ge m}(1+p_j)\,I=v$, $m$ the image's variables, RAM being zero above the image. The rows compute $\varphi^{-i}(\chi_j)=\chi_j^{2^{192-i}}$ from $\chi_j^{2^{128}}$, two Frobenius maps on the limbs, by squarings. A first-level node over a leaf circuit's proofs leaves, for each leaf, its hinted fixed-column evaluations, claims on $\fixleaf$, and its \textsf{HASH} matrix claim. A node's rows leave, for each child, its hinted fixed-column evaluations and its \textsf{HASH} matrix claim, and read the claims it carries off its statement, each weighed by whether the child's kind carries it.

The claims are then reduced by a small proof of its own, which the node's prover makes natively and the node's rows verify. Its transcript, seeded by a label digest and the zero statement, first binds every verified proof's terminal transcript commitment and every hint the claims rest on; each child's statement, carried claims included, is bound by its own transcript's seed. Then:
\begin{itemize}[itemsep=1pt,topsep=3pt]
\item \emph{The dense reduction.} With a challenge $\rbat$, the claims $s_t\,Q(p_t)=s_t\,v_t$ on the tree's dense polynomials $Q$, $s_t$ a claim's scale, become $\sum_Q\sum_x\dwt{Q}(x)\,Q(x)=\sum_t\rbat^t s_tv_t$ with $\dwt{Q}=\sum_{t\text{ on }Q}\rbat^t s_t\,\eq(p_t,\cdot)$, proven by one sumcheck of degree two over the most variables $n$, the lowest variable first. A polynomial of $n_Q<n$ variables is bound early and waits on the rest: its share is multiplied by each later challenge. Its end is checked against the prover's values $v_Q$, which become the claims $Q({\fdense}_{<n_Q})=v_Q$. The claims on $\fixnode$ or $\fixleaf$ from one child share their points' prefix, so the rows evaluate their $\eq$ factors on that prefix once.
\item \emph{The matrix reduction.} With a challenge, the claims $u^{\top}(aA_0+bB_0)w=v$, the fresh ones with $u=\erow$, $w=\wcol$ and $(a,b)=(1,\lcalpha)$ (\S\ref{flock:native}), the carried ones with $u=\eq(\frow,\cdot)$, $w=\eq(\fcol,\cdot)$ and $(a,b)$ one of $(1,0)$ and $(0,1)$, are batched and proven by a row phase over $\sum_f\sum_x\sum_c\rbat^cu_c(x)\,g_c(x)$ with $g_c=(a_cA_0+b_cB_0)w_c$, then a column phase over $\sum_f\sum_y A_0(\frow,y)\,\mwt{f}{A}(y)+B_0(\frow,y)\,\mwt{f}{B}(y)$ with $\mwt{f}{A}=\sum_c\rbat^cu_c(\frow)\,a_c\,w_c$: Spartan's second sumcheck~\cite{Set20}, every circuit sharing the challenges and waiting as above. The rows evaluate each $u$ and $w$ in closed form, $\erow$ through the skip domain's Lagrange weights and $\wcol$ through the slices. A leaf's fresh claims come from one batch (\S\ref{flock:batch}): they share $\zskip$, and those of one block size their whole row weight and the equality factor of their column weight, which the rows evaluate once, as they do each circuit's waiting factor. The prover computes $g_c$ by one forward walk of the circuit and $A_0(\frow,\cdot)$, $B_0(\frow,\cdot)$ by two backward walks (\S\ref{flock:matrices}), so the matrices are never built.
\end{itemize}
If an input claim is false, the batched claim is false except with probability at most $N/|\E|$ over $\rbat$, $N$ the number of claims, and a sumcheck of degree two over $n$ variables then ends on a false claim except with probability at most $2n/|\E|$. Every challenge is drawn after every input claim is bound. A claim whose scale is zero weighs nothing: a node scales each claim a child carries by the sum of the indicators of the kinds that carry it (one, and no scale, when every kind does), so a first-level node's carried claims on the polynomials its family does not fix, which carry no information, are not reduced.

\subsection{The root}\label{rec:root}

The root's verifier builds every kind's circuit, $\fixleaf$ and $\fixnode$ from the families and the shape, holds the root's kind to the one the expected tree's top gives, verifies the root's recursion proof against its statement, recomputes the digest from the expected tree, and settles the claims the root carries: each dense polynomial its kind carries a claim on (a first-level node's family's, every one for a node) at its prefix of $\fdense$, and each circuit's $A_0$ and $B_0$ at its prefixes of $\frow$ and $\fcol$, by one forward walk each. That is one evaluation of each polynomial whatever the number of leaves.

\paragraph{Soundness.} Suppose the root's verifier accepts. By the soundness of the root's recursion proof and \S\ref{rec:statement}, the root's rows hold on an assignment whose free wires are its children's statements and proofs. Each child's replay is then its verifier accepting its proof against its statement, except that its fixed columns' evaluations are hints and its \textsf{HASH} matrices are not evaluated: both are claims of the reduction's input. By the reduction's soundness, if any input claim is false, an output claim is false but with probability at most $(N+4\max_fk_f+2n)/|\E|$. The root's output claims are settled true, so every claim its children left or carried is true: each child's circuit is the one its kind names, its proof verifies, and the claims it carries are true. By induction down the tree every first-level node's rows hold on true claims, so each leanVM leaf's verifier's core accepts it and its deferred claims are settled, and each leaf circuit's proof verifies against its circuit, the one $\fixleaf$ names at the family's slot, so its statement holds of that circuit: every leaf verifies against the statement its first-level node's digest commits. The digests bind the statements, in order and with the tree's shape and the leaves' families, to the root's: a root over subtrees of different depths or families is provable, since a node verifies any kind, and its digest equals another tree's only by a collision of BLAKE2s, the levels' tags, the families' seeds and the counts keeping the messages apart. Knowledge soundness of the recursive composition of Fiat-Shamir proofs is, as for one level, heuristic in the random-oracle sense~\cite{BCCT13}.

\paragraph{Related work.} RISC Zero lifts each segment's proof into a recursion proof and joins them in pairs~\cite{BG23}; SP1 compresses its shard proofs by a tree of recursion proofs~\cite{SP1}; OpenVM's leaf circuit verifies a configurable number of application proofs and its internal circuit a configurable number of its own~\cite{OpenVM}. A first-level node here verifies the RISC-V proofs directly, as SP1's and OpenVM's first levels do, with no lift. Those systems bind the circuits they accept by a commitment to their verifying keys checked in the circuit. Here a circuit's identity is a claim on $\fixnode$ that every node reduces and only the root's verifier settles.
